% ================= 8. CONCLUSIONS & NEXT STEPS =================
\section{Conclusions \& Next Steps}
\label{sec:conclusions}

\subsection{Conclusions}
\label{sec:conclusions_summary}



The team has demonstrated the ability to build and control \projectname{} across the planned difficulty progression: a working 1-DoF model, and edge balance on the final 3D prototype, both achieved through the staged validation approach described in this document. Several of the originally planned tasks have been rescheduled in response to deadline pressure and issues encountered along the way, but the team has continued to make progress toward the final objective. Corner balance remains the outstanding milestone and is expected to be resolved in the days before final delivery and presentation.
\projectname{} has confirmed itself to be a tightly coupled, multidisciplinary system: reliably building and controlling it depends on close coordination and continuous communication across the team's subsystem interfaces.

\paragraph{Edge and corner balancing.} The outcome of both the edge and the corner test campaigns is judged satisfactory, because in each case the cube holds its unstable equilibrium with comfortable stabilisation margin and restores itself from displacements that meet or exceed what the acceptance criteria ask for. On the edge, balance is sustained unassisted past the Gate~G3 threshold and recovery was obtained from every angle of Table~\ref{tab:recovery_angle} and from the \qty{24}{\degree} hard stop of the rig support: requirement PER-03 is therefore satisfied with margin, and the demonstrated restoration envelope is bounded by the test fixture rather than by the controller (Section~\ref{sec:test_1dof}). At the corner, the nonlinear simulation campaign of Section~\ref{sec:per-corner} recovers at all eight resting corners, with a recoverable tilt of $2.76$--$3.14^\circ$ and a spread of only \qty{14}{\percent} between the best and worst corner once each is given its own gain set. The absolute corner envelope is small, but Table~\ref{tab:per-corner} shows why: every corner reaches the wheel-speed cap while the commanded torque stays inside its clamp, so the limit is the momentum budget of the wheels and not the quality of the stabilisation. Read together, the two campaigns say that the balancing law is not marginal at either equilibrium --- it recovers reliably wherever the actuators still have authority to act.

\paragraph{Control-loop refinement: auto-set and low-pass filtering.} Most of that margin was gained in the final iteration of the control loop, and two changes account for it. The first is the automatic equilibrium-set (\texttt{auto\_set}) mode, which establishes the balance reference of the current resting configuration online instead of relying on a hand-entered offset. The reference is not the same at every configuration --- the equilibrium tilt ranges from $0.714^\circ$ to $3.171^\circ$ across the eight corners of Table~\ref{tab:per-corner} --- and an incorrect reference is indistinguishable to the controller from a permanent disturbance: it is what produced the sustained oscillation observed with an uncalibrated centre of mass (Table~\ref{tab:balance_examples}) and what drains momentum budget continuously into a standing wheel speed (Table~\ref{tab:gyro-bias}). Setting the reference automatically removed that standing bias and, with it, the oscillation. The second is the low-pass filtering of the accelerometer path in the attitude estimator (Section~\ref{sec:cancellation}). The nonlinearity budget of Table~\ref{tab:nonlinearity-budget} identifies accelerometer noise as the dominant contributor to both tilt rms and torque activity, an order of magnitude above gyro noise, and the same path carries the lever-arm transient generated by the cube's own motion; attenuating it kept that noise out of the commanded torque without adding the phase lag that a heavier filter would have cost.

Both effects are largest at the corner. The corner has three coupled axes instead of one, so estimator noise is injected into three torque channels at once; its recovery envelope is momentum-limited, so any momentum spent chasing noise or a wrong equilibrium reference is taken directly out of the envelope; and its equilibrium tilt is both larger and more variable between corners than on the edge. This is precisely where the improvement was observed: the corner simulations became visibly quieter and their stabilisation markedly more reliable once the two changes were in place, and the corner results reported in this document are those of the refined loop.

\paragraph{Jump-up and walking.} The two stretch objectives, jump-up and walking (Section~\ref{sec:objectives}, Table~\ref{tab:traceability}), have not been demonstrated on the delivered cube, and the requirements that depend on them --- FUN-10, FUN-11, PER-07, PER-10 and CON-02 (Table~\ref{tab:requirements}) --- remain open in Appendix~\ref{app:verification}. The cause is not a gap in the underlying analysis: the jump-up sizing chain of Section~\ref{sec:jumpup_target} closes, the required wheel inertia is verified against it (PER-06, Table~\ref{tab:requirements}), and the mass and mechanical-interface budgets carry the brake servos throughout the design rather than as an afterthought --- they are costed in the mass budget (Table~\ref{tab:massbudget}) and given a controlled interface, ``Brake mount'', alongside the motor and battery interfaces (Table~\ref{tab:mech_interfaces}). What is missing is the mechanism itself: the servo bracket and the wheel strike feature, parts M5 and W4 of the CAD part list (Table~\ref{tab:cad_parts}, Appendix~\ref{app:cad}), were never carried from a reserved interface to a drawn and fabricated part, so no brake was ever mounted on the integrated cube and the manoeuvre was never attempted.

This gap is the direct consequence of a staging decision, taken at the outset and recorded in Section~\ref{sec:schedule}: balance the lighter, mechanically simpler cube first, and add the brake in a second design pass once that control problem was solved. The ordering was deliberate, not an oversight. A cube carrying three brake servos and their mounts is heavier and has a different mass and inertia distribution than one without, and an already-marginal balancing problem (Section~\ref{sec:background}) is harder to close, not easier, with that mass present from the first bring-up. Proving the balancing controller on the lighter configuration before committing to the brake's added mass was meant to let the team iterate quickly to a brake-equipped second version once the control loop was trusted, rather than debug an unstable equilibrium and a partially designed actuator at the same time. Consistent with this, the jump-up demonstration and the brake mechanism build were designated sacrificial scope from the start of the plan, sequenced after the edge-balance gate and positioned to yield first if time ran short (gate G8 of Table~\ref{tab:gates}; risk RE12 and cut-order item~2 of Table~\ref{tab:cascade}, Appendix~\ref{app:risks}).

That second pass never happened. Schedule pressure and the setbacks recorded in Appendix~\ref{app:risks} consumed the float the plan had reserved for it, and the team reached the final week with edge and corner balance --- the committed objectives of Section~\ref{sec:objectives} --- still the active critical path. At that point, integrating the brake meant reopening the electronics build on a heavier, structurally revised frame, and the drive, sensing and control electronics exist as a single set with no spare (RE11, Appendix~\ref{app:risks}): building a brake-equipped cube meant moving that one electronics set into a new prototype and re-establishing balance on it from scratch, with a real chance of reproducing, on the new build, the same stabilisation difficulties that had only just been solved on the current one. Weighed against the committed corner-balance and slew objectives, that risk was judged unacceptable, and the standing decision was to protect the current cube's balancing result rather than force a late brake integration under time pressure. This is a common and, in the team's judgement, correct trade in engineering practice: a verified baseline is worth more than an unverified stretch capability chased at its expense, and the jump-up and walking objectives are carried forward as the clearly scoped starting point of the next design iteration rather than rushed into this one. What that iteration consists of --- the brake mechanism itself, and the stiffer structure the braking impulse demands --- is set out in Section~\ref{sec:next_brake}.


\subsection{Next Steps}
\label{sec:next_steps}

The work that follows this issue falls into two streams, and they are not of the
same kind. The first is the completion of the cube itself: designing, building
and testing the brake mechanism, together with the reinforced structure that has
to survive it, so that the jump-up manoeuvre left open above can actually be
attempted. That stream is committed work on an existing design ---
\projectname{} is still in the process of being built, and the parts it needs
are already carried as reserved interfaces, costed in the mass budget and listed
as unbuilt items in the CAD part list. The second stream is exploratory: an
adaptive, learning-based outer loop wrapped around the verified controller,
motivated by the identification and tuning limits met during this project. It is
a research direction, not a scheduled deliverable, and nothing in the delivered
design depends on it. The two are kept apart below because they carry very
different levels of commitment and very different evidence behind them. Where
they meet --- Module~3, the optimisation of the jump-up torque profile --- the
learning work is explicitly downstream of the hardware: it cannot begin until
the brake exists and has been measured.

\subsubsection{Completing the Build: Brake Mechanism and Reinforced Structure}
\label{sec:next_brake}

\paragraph{The brake mechanism (parts M5 and W4).} The concept is already fixed
and does not need reopening: a servo-driven barrier is held in the wheel's path
and struck by a bolt head on the wheel, so the servo only \emph{holds} the
barrier while the wheel's own momentum does the work of the collision
(Section~\ref{sec:jumpup_target}). This is why the selected servo can be small
--- the TowerPro MG92B is rated near \qty{0.29}{\newton\metre}, about one
seventeenth of the \qty{\tauBrake}{\newton\metre} sizing torque --- and why the
design task is bounded rather than open-ended. What must still be produced is
the hardware itself: the brake servo bracket (M5, $\times 3$) and the wheel
strike feature (W4, $\times 3$), both still marked \emph{Not built} in
Table~\ref{tab:cad_parts}. The surrounding design already anticipates them: the
servos are costed in the mass budget (Table~\ref{tab:massbudget}) and a
controlled ``Brake mount'' interface is published alongside the motor and battery
interfaces (Table~\ref{tab:mech_interfaces}), so the second pass fills a
reserved slot instead of inventing one. The open geometry is the barrier travel
and its hard-stop position, the strike-face geometry on the wheel, and the bolt
pattern and stiffness of the bracket that reacts the impact.

\paragraph{A structure sized for the impact, not for balancing.} The cube as
built is a balancing structure. Every load-bearing part of it --- motor
subframe, motor mounts, wheel hub and spokes, frame members, corner contact
features --- is a 3D-printed part proportioned for the quasi-static torques of
stabilisation, where the largest force the frame ever sees is the reaction of a
commanded torque well inside its clamp. The braking event is a different
regime: it is a collision, and the two torque figures that describe it differ by
roughly an order of magnitude. The sizing value
$\tau_b = \qty{\tauBrake}{\newton\metre}$ is an impulse-equivalent mean,
$\bar\tau_b = h_w/\Delta t$, and implies a stopping time of about
\qty{52}{\milli\second}; a bolt head striking a printed barrier should arrest
the wheel in a few milliseconds, and the peak torque through the load path rises
in the same proportion (Section~\ref{sec:jumpup_target}). That load path is
specific and can be traced: barrier $\rightarrow$ bolt $\rightarrow$ servo
bracket $\rightarrow$ motor subframe on one side, and, in reaction, wheel rim
$\rightarrow$ spokes $\rightarrow$ hub $\rightarrow$ rotor adapter on the other,
with the resulting body impulse passing through the inner-structure-to-outer-frame
joint of Table~\ref{tab:mech_interfaces} and out through the corner contact
feature into the ground.

The next build is therefore not the present cube with a bracket added to it. It
is a structural redesign of the members on that path, checked against the
short-$\Delta t$ peak rather than against the sizing mean: thicker and ribbed
sections where the impulse is reacted, metal inserts or a metal bracket at the
strike and servo interfaces rather than a printed thread, print orientation
chosen so the impact loads the layers in compression and shear rather than
peeling them apart, and a stiffened frame joint so the impulse reaches the
ground without the outer frame flexing and stealing part of it. Stiffness
matters here for a dynamic reason as much as a strength one: energy absorbed in
a compliant frame during the collision is energy that does not go into tipping
the cube, so a structure that merely survives the brake is not sufficient --- it
has to be rigid enough to transmit the impulse, which is what makes the jump-up
succeed rather than merely not break anything.

The measurement that makes this design possible is the spin-down test already
scheduled as WBS task E10: it reports both the impulse-equivalent mean
$\bar\tau_b$ (the sizing quantity) and the peak torque (the structural one), and
its logging rate must resolve a few-millisecond event without smoothing the peak
away. Until it is run there is no defensible number to size the reinforced parts
to, which is why the measurement is the first item of this stream and not a
verification step at the end of it. The same test closes PER-10, which requires
the brake to arrest a wheel from the design speed within \qty{100}{\milli\second}.

\paragraph{What the rebuild costs the balancing loop.} Three brake servos, their
brackets, the strike features and the added structural material change $m$,
$r_{cm}$ and $I$ together --- the same three parameters whose uncertainty
motivates the adaptive work of Section~\ref{sec:ai_control}. The growth is
bounded and predicted rather than discovered, since the servos are already in
the mass budget and the total sits below the \qty{\Mcap}{\kilo\gram} cap of
CON-02, but the rebuilt cube is nonetheless a different plant: the
identification of Section~\ref{sec:paramid} has to be re-run on it, the
equilibrium references re-established, and the corner gain sets recomputed. Two
results of this project make that cheaper than it was the first time. The
\texttt{auto\_set} mode establishes the balance reference of each resting
configuration online, so a shifted centre of mass no longer has to be measured
by hand; and the per-corner tuning of Table~\ref{tab:per-corner} is now a
procedure that has been executed once rather than a problem to be solved from
scratch. The cost that does not fall is the electronics: drive, sensing and
control exist as a single set with no spare (RE11, Appendix~\ref{app:risks}), so
the transfer to a reinforced frame is a full re-build of the electrical
assembly and should be planned as one.

\paragraph{Demonstration sequence.} The order is set by containment --- each
step is run where a failure is cheap before it is run where it is not. First,
the spin-down measurement on a restrained single wheel, yielding $\bar\tau_b$
and the peak. Second, the structural check and the fabrication of the reinforced
parts against that peak. Third, a single-axis jump-up on the 1-DoF rig, where
the printed hard stops bound the excursion and the cube itself is not at risk.
Fourth, the integrated manoeuvre on the assembled cube: face to edge first, then
corner, with capture into the balancing law at the handover threshold. Fifth,
consecutive jumps, at a cadence no faster than roughly one every three seconds
given the \qty{2}{\second} spin-up of PER-12 plus the landing transient. Clearing
the fourth step closes FUN-10 and gate~G8 (Table~\ref{tab:gates}); the fifth
closes FUN-11; PER-07 and PER-10 close on the bench measurements that precede
them, and CON-02 is re-verified by weighing the reinforced build. All of these
remain \emph{Open} in Appendix~\ref{app:verification} at the date of this issue,
and closing them is the stated objective of the next iteration.

\subsubsection{AI-Driven Control \& Dynamic Coefficient Tuning}
\label{sec:ai_control}

\paragraph{Motivation: limitations of a fixed gain.} The balancing controller of Section~\ref{sec:lqr} computes a single
gain $K_d$ offline, from an $(A,B)$ pair built from the parameters of
Tables~\ref{tab:2d_vars} and \ref{tab:3d_vars}. That gain is optimal for
exactly one plant: the one those numbers describe. It is retained here as the
baseline precisely because it is simple and verifiable, but its optimality is
conditional in three ways that the hardware of this project is expected to
violate.

\begin{itemize}
    \item \textbf{Parametric drift.} $C_b$ and $C_w$ are not obtained from
    CAD and are carried as estimates until the identification tests of
    Table~\ref{tab:params}.
    Bearing friction,
    wiring drag and motor cogging also change with temperature, wear and
    reassembly, so the identified values are a snapshot rather than a
    constant. The inertia tensor $I$ and the centre-of-mass offset $r_{cm}$
    are subject to the same objection at the few-percent level, and $r_{cm}$
    in particular shifts whenever the pack or harness is disturbed.
    \item \textbf{Payload and configuration changes.} Any added mass changes
    $m$, $r_{cm}$ and $I$ simultaneously, and therefore changes $A$. A gain
    tuned for one configuration is not merely suboptimal for another; near a
    fully unstable equilibrium it may be destabilising.
    \item \textbf{Contact and surface variation.} The pivot is modelled as an
    ideal frictionless point or line. A real corner resting on a real surface
    slips, deforms and dissipates energy in a way that is neither in the model
    nor constant across surfaces, and this is the dominant unmodelled effect
    for the walking objective of Section~\ref{sec:objectives}.
\end{itemize}

The standard remedy within the existing design is gain scheduling on
$\hat\theta_b$ (Section~\ref{sec:lqr}), but scheduling interpolates
between gains that were themselves computed from the same assumed parameters.
It addresses the weakening of the small-angle linearisation; it does not
address a plant that differs from the model. The proposal here is to keep the
verified feedback law and adapt its coefficients online.

\paragraph{Proposed hybrid architecture.} The intended structure is deliberately conservative: a classical inner loop
whose stability properties are understood, wrapped by a slow outer loop that
adjusts the inner loop's coefficients. Writing $\phi$ for the tunable
coefficient vector,
\begin{equation}
u[k] = -K_d(\phi[k])\,\hat x[k],
\qquad
\phi[k+1] = \phi[k] + \Delta\phi\big(\hat x[k],\, \hat x[k-1], \dots\big),
\label{eq:hybrid_arch}
\end{equation}
where $\hat x$ is the estimated state of Section~\ref{sec:lqr} and
$\Delta\phi$ is produced by the tuner. Two properties of
\eqref{eq:hybrid_arch} matter more than the choice of learning algorithm.
First, the inner loop remains a linear state feedback at every instant, so
the controller that runs on the hardware is the one already implemented and
tested. Second, $\phi$ is updated on a timescale far slower than the control
period $\Delta t$, which keeps the frozen-gain analysis approximately valid
and prevents the tuner from acting as an unmodelled fast feedback path. Three
candidate tuners are of interest, and they are not equivalent in cost:

\begin{itemize}
    \item \textbf{Bayesian optimisation} over a low-dimensional
    parameterization of $Q$ and $R$, minimising a measured cost from short
    balancing trials \cite{shahriari2016bayesopt}. This is the cheapest option
    and the one with a direct precedent on this class of hardware: automatic
    LQR tuning has been demonstrated on an inverted pendulum with on the order
    of tens of experiments \cite{marco2016automaticlqr}, and a self-tuning LQR
    has been run on the Cubli's own ancestor platform at ETH
    \cite{trimpe2014learningcubli}. It tunes between trials, not within them.
    \item \textbf{Adaptive dynamic programming}, which solves the LQR problem
    online from measured data rather than from $(A,B)$, converging toward the
    optimal gain for the true plant without an explicit model
    \cite{lewis2012adp}. This is the option with the strongest theoretical
    guarantees and the strictest excitation requirements.
    \item \textbf{Reinforcement learning}, in which a policy trained in
    simulation outputs $\Delta\phi$ continuously from the state-error history.
    This is the most capable and the least verifiable, and is treated below as
    the exploratory branch rather than the default.
\end{itemize}

\paragraph{Deployment reservations.} The literature supporting learned control
is overwhelmingly validated in simulation or on well-instrumented hardware,
and the failure mode of an online tuner on an open-loop-unstable plant is a
fall, not a degraded score. Any deployment on this cube should therefore be
gated behind: a projection of $\phi$ onto a set for which the frozen-gain
closed loop is verified stable across the parameter uncertainty; a fallback to
the fixed $K_d$ of Section~\ref{sec:lqr} on any estimator fault or
divergence; and bench trials with the cube tethered. This is future work in
the literal sense --- none of it is a committed deliverable of the present
design, and the baseline controller is not contingent on any of it.

\paragraph{Module 1 --- dynamic gain tuning.} The tuner does not emit motor currents. It emits the coefficients from which
$K_d$ is computed, which bounds what a bad update can do. Two
parameterizations are worth comparing:

\begin{itemize}
    \item \textbf{Weight-space tuning.} The policy outputs the diagonal
    entries of the cost matrices,
    $\phi = (\log q_1,\dots,\log q_9,\ \log r_1, \log r_2, \log r_3)$ with
    $Q = \mathrm{diag}(q_i) \succ 0$ and $R = \mathrm{diag}(r_i) \succ 0$, and
    $K_d$ follows from the discrete algebraic Riccati equation. The
    logarithmic parameterization enforces positive definiteness by
    construction, so every $\phi$ the policy can emit corresponds to a
    well-posed LQR problem and, for a stabilisable $(A_d,B_d)$, to a
    stabilising gain for the \emph{modelled} plant. The cost is that a Riccati
    solve must run onboard whenever $\phi$ changes, which is what confines
    updates to a slow outer loop.
    \item \textbf{Gain-space tuning.} The policy perturbs the entries of $K_d$
    directly, or equivalently the $K_p, K_d$ pair of a PD form, avoiding the
    Riccati solve at the price of losing the structural guarantee above. This
    is the practical choice if the outer loop must run fast, and it makes the
    projection step of the previous paragraph mandatory rather than merely
    prudent.
\end{itemize}

For the 3D corner-balancing case the state is
$x = (\theta, \omega_b, \omega_w) \in \mathbb{R}^9$ and $K_d$ is $3\times 9$
(Section~\ref{sec:3dmodel}), so tuning the $54$ free entries directly is not
sensible on the available data. Weight-space tuning reduces the problem to
$12$ coefficients, and symmetry of the three wheel axes reduces it further if
the wheels are treated as identical. A reward or cost of the form
\begin{equation}
c[k] = \hat x^{\mathrm T}[k]\,W_x\,\hat x[k] + \rho\,\|u[k]\|^2 + \sigma\,\|\omega_w[k]\|^2
\label{eq:tuner_cost}
\end{equation}
Equation~\eqref{eq:tuner_cost} adds an explicit penalty on stored wheel momentum, which the nominal LQR cost
$J(u)$ of Section~\ref{sec:lqr} does not distinguish from any other
state deviation. This term is the one directly relevant to the saturation
problem raised in Section~\ref{sec:application}: it expresses a preference for
holding attitude with the wheels near zero speed, leaving momentum headroom
for disturbance rejection.

\paragraph{Module 2 --- sim-to-real transfer via domain randomization.} A policy trained against a single nominal parameter set learns to exploit that
parameter set. Domain randomization instead samples the plant parameters from
a distribution at the start of each training episode
\cite{tobin2017domainrand}, so that the policy is optimised against a family
of plants rather than one, and the transfer method has a substantial record on
physical legged hardware \cite{hwangbo2019anymal}. For this cube the
randomized quantities follow directly from the identification gaps already
documented:

\begin{itemize}
    \item \textbf{Damping:} $C_b, C_w$ over a wide interval, since these are
    estimated rather than measured until the identification of
    Section~\ref{sec:paramid} closes.
    \item \textbf{Inertia and mass distribution:} $I$, $m$ and $r_{cm}$
    perturbed by a few percent, covering print-density variation, fastener
    placement and harness routing.
    \item \textbf{Actuation:} $K_m$ within motor tolerance, plus a current
    saturation limit and a commanded-torque delay, so the policy cannot assume
    an ideal actuator.
    \item \textbf{Contact:} pivot friction and restitution, and a small random
    ground slope --- the term standing in for the surface variation that
    motivates this module.
    \item \textbf{Sensing:} gyro bias walk, accelerometer noise and an
    injected vibration component, matching the estimator degradation
    identified in Section~\ref{sec:background} as the practical limit on
    balancing.
\end{itemize}

The sensing term requires particular care. The accelerometer corruption on this
platform is generated by the controller's own wheel activity, so it is
correlated with the policy's actions rather than independent of them. A
randomization scheme that injects only white sensor noise will understate the
difficulty, and a policy that transfers under it should not be assumed to
transfer on hardware. Whether the randomized distribution covers the real cube
is an empirical question, settled by bench testing; the principal failure mode
of this module is a distribution selected for convenience rather than
representativeness.

\paragraph{Module 3 --- jump-up energy optimisation.} This module presupposes the
hardware of Section~\ref{sec:next_brake}: it optimises the use of a brake that
must first be built and measured. The jump-up analysis of Section~\ref{sec:jumpup_target} treats the manoeuvre as
a spin-up to $\omega_{\max}$ followed by an abrupt stop, with all deviation
from the ideal absorbed into a single brake-amplification factor
$\beta = \tau_b/(\tau_b-\tau_g)$. That model is adequate for sizing and is
deliberately conservative, but it is a poor description of what the mechanism
does: the servo-driven barrier has finite travel, the collision is not
instantaneous, and the resulting torque profile is neither a step nor known in
advance.

The proposal is to treat the torque command over the manoeuvre window as a
profile to be optimised rather than a constant. Parameterizing the
pre-brake and post-impact torque commands by a small vector $\psi$ --- for
instance the timing of brake release relative to the wheel-speed trajectory,
and the shape of the recovery torque applied by the remaining wheels
immediately after impact --- the objective is
\begin{equation}
\min_{\psi}\ \ \mathbb{E}\Big[\, \|\theta(t_f) - \theta^\star\|^2 + \mu\,\|\omega_b(t_f)\|^2 \,\Big]
\quad \text{subject to} \quad
\omega_w \le \omega_{\max},\ \ |u_i| \le u_{\max},
\label{eq:jump_opt}
\end{equation}
where $t_f$ is the instant the controller hands over to the balancing law at
$|\hat\theta_b| < \theta_{\text{th}}$. Minimising residual body rate at
handover is the quantity of interest, because a jump that reaches the
equilibrium with excess rate demands recovery authority the balancing
controller may not have. Because each evaluation of \eqref{eq:jump_opt} is a
physical jump, the sample budget is small and the appropriate method is
Bayesian optimisation rather than a policy-gradient method, with the initial
design taken from the analytical $h_w$ of Equation~\eqref{eq:hw_ideal}.

Two constraints on this module must be stated. First, it presupposes the
measurement described in Section~\ref{sec:jumpup_target}: until the spin-down
test reports both the impulse-equivalent mean
$\bar\tau_b = h_w/\Delta t$ and the peak torque, there is no validated model to
optimise against and no way to distinguish an improved profile from a
mis-measured one. Second, the structural caveat carries over unchanged ---
optimisation must not be permitted to search toward shorter stopping times,
since $\tau_b$ is simultaneously a load on the barrier, bolt, bracket and
spokes, and the sizing value was chosen for the dynamics and not for that load
path. The search space must be bounded by the structural case, not by the
dynamic one.

\paragraph{Relevance beyond the bench.} Two application claims follow, and both are narrower than the general
enthusiasm for learned control would suggest.

For \textbf{CubeSat attitude control}, the relevant fact is that a
spacecraft's inertia tensor is known imprecisely at integration and changes in
flight --- propellant depletion, deployable appendages, thermal distortion ---
while reaction-wheel friction rises with bearing age. These are the orbital
analogues of the parametric drift of Section~\ref{sec:lqr}, and they are
conventionally handled by conservative fixed gains with margin, at a cost in
pointing performance. A tuner that adapts $Q$ and $R$ against measured
performance addresses exactly this gap. The qualification barrier for flying
learned control is high and unaddressed here; the architecture of
\eqref{eq:hybrid_arch} is chosen partly because a bounded outer loop over a
classical inner loop is the form most likely to survive that scrutiny.

For \textbf{low-gravity surface mobility}, the case is stronger, because the
plant is genuinely unknown before arrival. Regolith mechanical properties on a
small body are not characterised to the precision a hopping controller would
want, and the flight-proven hoppers cited in Section~\ref{sec:application}
\cite{yoshimitsu1999minerva, ho2017mascot} move ballistically rather than
under closed-loop control for this reason. The contact randomization of
Module~2 and the profile optimisation of Module~3 map onto that problem
directly: a hop whose landing attitude is controlled, on a surface whose
properties were sampled from a distribution rather than assumed, is the
capability that separates directional walking from a ballistic tumble
\cite{spiridonov2024spacehopper}. The cube tests this under $1$~g, where the
destabilising torque is larger and never relents --- a harder case in that one
respect, as argued in Section~\ref{sec:application}, though it does not
reproduce the low-gravity contact dynamics that would decide the real
manoeuvre.

% \subsection{Single-Wheel Variant / Extensions}

% A reduced-actuation variant balancing the 3D cube with a single reaction wheel,
% following the ETH one-wheel Cubli \cite{hofer2023onewheelcubli}.
